\problemname{Santa Planning Building}
According to the UN, the world's population reached 8 billion people in November 2022. One person who clearly
notices this is Santa Claus. He has to deliver presents to several billion children.

Unfortunately, Santa will not have time to deliver presents to all the children this year. Instead, he needs
to choose a group of children to receive presents. To make it fair, he plans to give presents to exactly half
of all the children, while the other half will go without. This way, the two groups will contain exactly the
same number of children, and no group will be envious of the other. ``Perfectly balanced'', as Santa often says.

You are responsible for planning which houses Santa should visit. There are $N$ houses numbered from $1$ to $N$,
where house number $i$ has coordinates $(x_i, y_i)$. Traveling between two points $(x_i, y_i)$ and $(x_j, y_j)$
takes $|x_i - x_j| + |y_i - y_j|$ seconds\footnote{$|a|$ is the mathematical notation for the \emph{absolute value}
of the number $a$, defined as $a$ if $a \ge 0$ and $-a$ if $a < 0$.}.
Your task is to find an ordered list of $N/2$ different houses ($N$ is always even), such that the total time to
visit the houses in that order is as small as possible. Santa starts his journey at the first house in your list
and finishes it at the last house in your list.

\section*{Input}
The input consists of $10$ test cases.

The first line contains an integer $T$ ($0 \leq T \leq 10$), the number of the test case ($0$ is the sample case below).

The second line contains an integer $N$ ($1 \leq N \leq 10^5$), the number of houses. $N$ is guaranteed to be even.

The following $N$ lines each contain two integers $x_i$ and $y_i$ ($0 \leq x_i, y_i \leq 10^6$), the coordinates of the
$i$-th house. No two houses are located at the same coordinates.

\section*{Output}
Print $\frac{N}{2}$ space-separated integers, the numbers of the houses to be visited in order.

\section*{Points}
Your program will be tested on $10$ test cases and scored based on how it compares to the judges' solution.

Let $P_1$ be the total distance Santa needs to travel in your solution, and let $P_2$ be the corresponding
distance for the judges' solution. You will then receive $\min(10 \cdot \frac{P_2}{P_1}, 10)$ points for the test case.
You can never receive more than $10$ points for a test case, even if your solution is better than the judges' solution.
However, the judge might show a higher score per test case in the submission details interface so that you can see exactly
how good your solution is.

Your score for a submission is the sum of the scores for the test cases, and your total score is the score of your best
submission. Thus, it is not possible to combine test cases from different submissions; instead, you must submit a
single program that maximizes your score on all test cases simultaneously.

The example test case will always give $0$ points.

\section*{Testcases}
Here are descriptions of the $10$ test cases that occur. In all test cases, the coordinates are guaranteed to have been
generated uniformly at random between $0$ and $10^6$, unless otherwise specified.

\noindent
\begin{tabular}{| l | l | p{12cm} |}
  \hline
  \textbf{Group} & \textbf{Point value} & \textbf{Constraints} \\ \hline
  $1$       & $N = 10^5$                                 & $y_i = 0$ for all $i$  \\ \hline
  $2$       & $N = 10^5$                                 & $0 \leq y_i \leq 20$ for all $i$   \\ \hline
  $3$       & $N = 10$                                   &    \\ \hline
  $4$       & $N = 40$                                   &   \\ \hline
  $5$       & $N = 200$                                   &    \\ \hline
  $6$       & $N = 800$                                   &    \\ \hline
  $7$       & $N = 7000$                                   &    \\ \hline
  $8$       & $N = 40000$                                   &   \\ \hline
  $9$       & $N = 80000$                                   &    \\ \hline
  $10$      & $N = 10^5$                                   &    \\ \hline
\end{tabular}

\section*{Explanation of sample 1}
In the example case, the example solution achieves a total distance of $40$ because the distance between
house $2$ and house $4$ is $5$, and the distance between house $4$ and house $1$ is $35$. There are several
ways to achieve a better answer, for example by visiting the houses in the order \verb|4 2 1|. Then the total
distance becomes $35$ instead.
